% =============================================================================
% Title: The Illusion of a Flat Earth within a 3-Torus Universe
% Author: Brendon Boyd
% Purpose: Zenodo Preprint (Bridge Paper: Academic Rigor + Public Accessibility)
% =============================================================================

\documentclass[aps,prd,11pt,onecolumn,superscriptaddress,nofootinbib,floatfix]{revtex4-2}

\usepackage[utf8]{inputenc}
\usepackage[T1]{fontenc}

% Mathematics
\usepackage{amsmath, amssymb, amsthm}
\usepackage{physics}

% Tables & Formatting
\usepackage{booktabs}
\usepackage{setspace}
\onehalfspacing

% References and links
\usepackage{natbib}
\usepackage[colorlinks=true, citecolor=blue, urlcolor=blue, linkcolor=black]{hyperref}

\begin{document}

\title{Locally Flat, Globally Periodic: \\ Reconciling the Illusion of a Flat Horizon via $T^3$ Topology}
\author{Brendon Boyd}
\email{brendon.boyd@itsm-cosmology.org}
\thanks{ORCID: 0009-0007-4177-2612 | Licensed under Creative Commons Attribution 4.0 International. ``ITSM Cosmology'' and ``Integrated Toroidal-Syntropic Model''\texttrademark\ are unregistered trademarks of Brendon Boyd. GitHub: \url{https://github.com/brendohxd/ITSM-Integrated-Toroidal-Syntropic-Model}. Zenodo: \href{https://doi.org/10.5281/zenodo.20773262}{10.5281/zenodo.20773262}}
\affiliation{Independent Researcher, Burswood, Western Australia, Australia}
\date{Compiled \today}

\begin{abstract}
Observational anomalies such as an apparent flat horizon, circular celestial paths, and perceived boundedness have driven a resurgence in flat earth theories. While the scientific community widely dismisses these conclusions as conspiratorial, the underlying geometric observations themselves warrant rigorous analysis. This paper proposes a novel resolution: these phenomena are not evidence of a flat planar Earth, but rather low-dimensional projections of a universe possessing a 3-torus ($T^3$) topology. Because a $T^3$ manifold is locally Euclidean (flat) but globally periodic, it mathematically accommodates a flat-appearing horizon while enforcing repeating, cyclical celestial boundaries. When coupled with the Integrated Toroidal-Syntropic Model (ITSM), this framework explains the "dome" illusion through superfluid shear planes and provides testable predictions, including topological lensing and repeating galaxy clusters observable by the James Webb Space Telescope (JWST). This model serves as a scientific bridge, providing fundamental topological mechanics to explain controversial celestial observations without violating modern astrophysics.
\end{abstract}

\maketitle


% =============================================================================
\section{Introduction}
\label{sec:intro}

In recent years, a segment of the public has challenged the prevailing cosmological model by pointing to direct, localized celestial observations—most notably, the perception of a perfectly flat horizon, the strictly circular paths of the sun and moon, and the illusion of a bounded "dome" or firmament. The standard scientific response has been to dismiss these observations as misunderstandings of scale and gravity on a spherical planet. However, dismissing the observations outright alienates highly analytical and inquisitive minds, creating an intellectual stalemate.

When independent researchers observe these geometric anomalies, their analytical deduction is often sound. But without access to advanced topological models, they are forced into restrictive conclusions—such as a planar Earth bounded by an infinite ice wall. A core motivation of this paper is to provide a rigorous mathematical bridge for these analytical minds. By equipping them with the framework of a 3-torus ($T^3$) topology, this model aims to redirect their observational acuity away from marginalized theories and empower them to actively contribute to the progress of modern astrophysics.

This paper proposes an alternative: validating the geometric integrity of the observations while correcting the conclusion. The phenomena observed do not require a planar Earth; rather, they are perfectly consistent with the mathematical properties of a $T^3$ universe. 

By analyzing local flatness, periodic boundary conditions, and topological lensing within a $T^3$ framework—supported by the dynamics of the Integrated Toroidal-Syntropic Model (ITSM)—we can mathematically explain these "flat earth" anomalies as low-dimensional projections of a highly structured universe.

% =============================================================================
\section{Mathematical and Physical Properties of the 3-Torus ($T^3$)}
\label{sec:t3-properties}

A 3-torus is the Cartesian product of three circles, $S^1 \times S^1 \times S^1$, forming a 3-dimensional manifold with periodic boundary conditions.

\subsection{Local Flatness vs. Global Periodicity}
The defining feature of a $T^3$ topology is that its intrinsic geometry is exactly Euclidean (flat). To a localized observer, a $T^3$ universe appears entirely flat, possessing zero spatial curvature. This is in perfect alignment with modern cosmological data, such as the Planck 2018 results, which measure the curvature parameter $\Omega_k \approx 0$ \cite{planck_2018}. 

However, unlike an infinite Euclidean space, the $T^3$ is globally periodic. Traveling far enough in any direction eventually returns the traveler to their starting point, much like navigating a 3D video game world. This periodicity means the universe has no edge, no center, and is finite in volume, while remaining spatially flat.

\subsection{Topological Lensing}
Because light can wrap around the periodic boundaries of a $T^3$ universe, an observer can theoretically see multiple images of the same celestial object originating from different directions and epochs. This topological lensing creates a cosmic "hall of mirrors" effect, potentially generating repeating patterns in the sky.

% =============================================================================
\section{Mapping Observations to $T^3$ Topology}
\label{sec:mapping}

When the observational claims of flat earth theorists are analyzed through the lens of a $T^3$ topology, the apparent contradictions with modern astrophysics evaporate. Table~\ref{tab:mapping} details this alignment.

\begin{table}[h]
\centering
\caption{Celestial Observations Mapped to $T^3$ Explanations}
\label{tab:mapping}
\begin{tabular}{p{0.25\textwidth} p{0.35\textwidth} p{0.3\textwidth}}
\toprule
\textbf{Observation/Claim} & \textbf{$T^3$ Topological Explanation} & \textbf{Astrophysical Context} \\
\midrule
Horizon appears flat & Local intrinsic flatness of the $T^3$ manifold. & Planck 2018 ($\Omega_k \approx 0$) \cite{planck_2018}. \\
Circular celestial paths & Geodesics on $T^3$ appear circular due to periodic boundaries. & Analogous to plasma confinement paths in tokamak reactors \cite{iter_basis}. \\
"Dome" illusion & Shear planes within the cosmic fluid create 2D-like boundedness. & Modeled by ITSM's superfluid plenum dynamics. \\
No observable curvature & $T^3$ is globally periodic, not curved; it is the simplest compact flat manifold. & Topological mathematics. \\
Celestial bodies "reset" & Periodic boundary conditions cause "wrapping" of light and object paths. & Topological lensing. \\
\bottomrule
\end{tabular}
\end{table}

\subsection{The "Dome" and Shear Planes}
One of the most persistent illusions in alternative cosmologies is the existence of a firmament or dome. In the context of the Integrated Toroidal-Syntropic Model (ITSM), the universe is filled with a superfluid plenum. Within this plenum, differential fluid motion creates 2D-like surfaces known as \textit{shear planes}. These topological defects act as localized boundaries or constraints on celestial mechanics, providing a rigorous fluid-dynamics explanation for the perception of a bounded, dome-like sky.

% =============================================================================
\section{Falsifiability and Predictions}
\label{sec:predictions}

A robust scientific model must offer testable predictions. The $T^3$ universe model presents several unique observational signatures that differentiate it from the standard $\Lambda$CDM model:

\begin{enumerate}
    \item \textbf{Repeating Galaxy Patterns:} Deep-field surveys, particularly those conducted by the James Webb Space Telescope (JWST), should reveal statistically significant repeating galaxy clusters in opposite directions of the sky, indicating light that has wrapped around the universe's periodic boundaries.
    \item \textbf{Topological Lensing Anomalies:} Gravitational lensing observations should reveal specific anomalous light paths and multiple imaging that cannot be accounted for by the mass of intervening foreground clusters alone.
    \item \textbf{Shear Plane Signatures in Rotation Curves:} The influence of superfluid shear planes should manifest as 2D-like motion in specific celestial bodies, potentially explaining anomalies in the SPARC galactic rotation curve database \cite{sparc_2016} without the need for cold dark matter.
    \item \textbf{Toroidal Casimir Anomalies:} Laboratory experiments measuring the Casimir effect \cite{casimir_1948} within highly constrained toroidal cavities should reveal measurable deviations from standard Euclidean predictions, reflecting the macroscopic topology at a quantum level.
\end{enumerate}

% =============================================================================
\section{Conclusion}
\label{sec:conclusion}

The dismissal of anomalous celestial observations as mere conspiracy theories stifles scientific inquiry. By taking the geometric observations of a flat horizon and cyclical celestial paths seriously, we find that they point not to a planar Earth, but to a beautifully complex, locally flat, and globally periodic 3-torus ($T^3$) universe. 

This framework liberates inquisitive minds from the restrictive binary of "spherical earth vs. ice wall" by introducing a third, mathematically rigorous paradigm. As deep-space observatories like JWST continue to probe the furthest reaches of the cosmos, the topological signatures of a $T^3$ universe may soon move from theoretical resolution to observational fact.

% =============================================================================
\section*{Acknowledgments}
The foundations of this topological approach are deeply intertwined with the development of the Integrated Toroidal-Syntropic Model (ITSM) \cite{itsm_main}. 

% =============================================================================

\begin{thebibliography}{9}
\bibitem{itsm_main} Boyd, B. (2024). \textit{Integrated Toroidal-Syntropic Model (ITSM)}. Zenodo. \href{https://doi.org/10.5281/zenodo.20773262}{doi:10.5281/zenodo.20773262}
\bibitem{planck_2018} Planck Collaboration (2020). \textit{Planck 2018 results. VI. Cosmological parameters}. Astronomy \& Astrophysics 641, A6. \href{https://doi.org/10.1051/0004-6361/201833910}{doi:10.1051/0004-6361/201833910}
\bibitem{iter_basis} ITER Physics Expert Groups et al. (1999). \textit{ITER Physics Basis}. Nuclear Fusion 39, 2137. \href{https://doi.org/10.1088/0029-5515/39/12/308}{doi:10.1088/0029-5515/39/12/308}
\bibitem{sparc_2016} Lelli, F., McGaugh, S. S., \& Schombert, J. M. (2016). \textit{SPARC: Mass Models for 175 Disk Galaxies with Spitzer Photometry and Accurate Rotation Curves}. The Astronomical Journal, 152(6), 157. \href{https://doi.org/10.3847/0004-6256/152/6/157}{doi:10.3847/0004-6256/152/6/157}
\bibitem{casimir_1948} Casimir, H. B. G. (1948). \textit{On the attraction between two perfectly conducting plates}. Proceedings of the Koninklijke Nederlandse Akademie van Wetenschappen, 51, 793–795.
\end{thebibliography}

\end{document}
